\section{Schema Migration}
\label{sec:schema-migration}

Phase 2 extends the frozen Phase 1 schema with maintenance impact levels,
advisory acknowledgements, and data-driven instant booking. Since the database
already contained booking and maintenance data, the upgrade was implemented as
a data-preserving T-SQL delta in \texttt{10-schema-migration.sql}. The script
upgrades the Task 05 baseline, \texttt{05-db-definition.sql}, to the approved
Task 09 design. The companion script
\texttt{10-schema-migration-rollback.sql} reverses the upgrade when required.

\subsection{Migration Approach}

The migration is primarily additive: it introduces new columns, tables,
indexes, and triggers without rebuilding or renaming any Phase 1 table. The
only removed element is \texttt{spaces.usage\_policy}, an unused free-text
column superseded by the Task 09 reference-table design. No rule, index, or
trigger depended on this column.

The complete schema delta is summarized in Table~\ref{tab:migration-delta}.

\begin{table}[htbp]
    \centering
    \small
    \renewcommand{\arraystretch}{1.15}
    \begin{tabularx}{\textwidth}{>{\RaggedRight\arraybackslash}p{0.27\textwidth}
                                  >{\RaggedRight\arraybackslash}X}
        \toprule
        \textbf{Schema component} & \textbf{Phase 2 implementation} \\
        \midrule
        Maintenance impact &
        Add \texttt{maintenance.impact\_level} as \texttt{NOT NULL}, defaulting
        to \texttt{'out-of-service'}. A check constraint permits only
        \texttt{out-of-service} and \texttt{advisory}. \\
        \midrule
        Impact history &
        Add \texttt{maintenance\_impact\_history} to record every maintenance
        impact-level change. \\
        \midrule
        Advisory acknowledgement &
        Add \texttt{booking\_advisory\_acknowledgement}. A composite unique
        constraint permits one acknowledgement per booking--advisory pair. \\
        \midrule
        Instant-booking policy &
        Add \texttt{space\_type\_allowed\_purpose} and remove the superseded
        \texttt{spaces.usage\_policy} column. \\
        \midrule
        Trigger logic &
        Restrict booking blocks to out-of-service maintenance; record impact
        history; validate acknowledgement ownership; maintain modification
        timestamps; and recompute space status. \\
        \midrule
        Configuration data &
        Seed one reserved system approver and 18 allowed-purpose pairs for the
        four approved space types. \\
        \bottomrule
    \end{tabularx}
    \caption{Phase 2 schema migration delta}
    \label{tab:migration-delta}
\end{table}
\FloatBarrier

The revised booking and maintenance triggers block overlaps only when a space
has out-of-service maintenance; advisory maintenance produces a warning but
does not make the space unavailable. Additional triggers record impact-level
changes, verify that acknowledgements are submitted by the booking owner,
maintain last-modified timestamps, and recompute
\texttt{spaces.current\_status}. Status is derived using the approved priority
order, avoiding duplicated state.

The migration is both transactional and repeatable. All operations execute in
a single transaction, so a failed or interrupted run leaves the database
unchanged. Existence checks guard object creation and seed insertion, allowing
the script to be rerun without creating duplicate objects or rows. The rollback
script restores the four Phase 1 trigger definitions and removes every column,
table, index, trigger, and seed row introduced by the migration, returning the
schema to the Task 05 baseline.

\subsection{Data Preservation and Backfill}

No existing booking or maintenance row is deleted or assigned a new identity.
The only backfill sets the new maintenance impact level to
\texttt{`out-of-service'}. This preserves Phase 1 behavior because every
existing active maintenance record continues to block overlapping bookings.

After the schema changes, \texttt{spaces.current\_status} is recomputed using
the priority order: retired, temporarily closed, out-of-service maintenance,
live session, and available. Existing manual \texttt{retired} and
\texttt{temporarily\_closed} overrides are preserved.

The migration also inserts the reserved system approver used to identify
derived instant approvals and the 18 \texttt{(space\_type, purpose)} pairs that
define instant-booking eligibility for the four approved space types. Guarded
inserts ensure that each row is created only once and do not disturb existing
identity values.

\subsection{Verification}

Verification was performed on a scratch SQL Server 2022 database. The Task 05
DDL and sample data were loaded first, followed by the migration, a second run
of the same migration, and the rollback. Validation queries confirmed that the
new schema objects had the approved definitions, every maintenance row had an
impact level, and the system approver and all 18 allowed-purpose pairs existed
exactly once. The second migration run completed without errors or duplicates,
confirming repeatability.

Transactional smoke tests covered the main behavioral rules. Advisory
maintenance allowed overlapping bookings, whereas out-of-service maintenance
continued to block them. Duplicate acknowledgement of the same advisory was
rejected, as was approval without the required acknowledgement. Each test was
rolled back after execution. Finally, the rollback script restored the
baseline row counts exactly, confirming that the migration and reversal caused
no data loss.